% Chapter introduction: why the topic arose, who first asked, what it answers.
% STATUS: written 2026-11-05 (all five body sections sourced and gate-clean).
% The introduction reports the four threads of the chapter and cites only the
% sources that carry them (Ferrero-Haetinger 2002, Isaias-Narvaez Macarro 2009,
% Vukman 1989 and 1990, Hongan-Komatsu 2000, Yazarli 2017, Prajapati 2021,
% Wang 2016, Liu-Liu 2021, Hasse-Schmidt 1937, Nakai 1970).
\emph{Why a single derivation is not enough.} Chapter~\ref{ch:generalized-twisted}
changed one map: from a derivation to a $\sigma$-derivation, an
$\alpha$-derivation, a semiderivation or a left derivation
\cite[p.~2841]{Chuang2017Auto}. This chapter changes the map in four more
ways, and each way is read from the paper that carries it.

\hbfigure{07-higher-multi/hasse-schmidt-grid}{The product rule of a higher derivation: $D_n(xy)$ collects the terms $D_i(x)D_j(y)$ on the antidiagonal $i+j=n$, here $n=3$.}{fig:higher:grid}

The first way replaces the derivation by a sequence
\cite[p.~2323]{FerreroHaetinger2002}. Hasse and Schmidt built the first
sequences in 1937, and their two rules hold for an arbitrary integral domain
\cite[p.~224]{HasseSchmidt1937}. Section~\ref{sec:higher-multi:hasse-schmidt}
recalls those rules and reads the definition from Ferrero and Haetinger
\cite[p.~2323]{FerreroHaetinger2002}. The single-map notion of a $q$-th order
derivation is Nakai \cite[p.~1]{Nakai1970}, and the extension theorems are
Narv\'aez Macarro \cite[pp.~2983--2984]{NarvaezMacarro2009}.

The second way keeps one map and gives it two arguments
\cite[p.~245]{Vukman1989Bid}. Maksa introduced the symmetric biderivation
\cite[references~5 and~6]{Vukman1989Bid} and Vukman carried it into prime
rings \cite[p.~245]{Vukman1989Bid}; Section~\ref{sec:higher-multi:biderivations}
recalls his Theorem~1 and Theorem~2 \cite[pp.~246, 249]{Vukman1989Bid}.

The third way lets the number of arguments grow
\cite[p.~490]{Yazarli2017Trideriv}. Section~\ref{sec:higher-multi:n-derivations}
introduces the permuting triderivation of Öztürk
\cite[reference~10]{Ozturk1999Trideriv} and runs the second theorem of Posner
along higher derivations \cite[p.~9]{Prajapati2021}.

The fourth way studies the trace on its own
\cite[p.~181]{Vukman1990Two} and asks what a biderivation can be on a
structured ring \cite[p.~1952]{Wang2016Bideriv}.
Sections~\ref{sec:higher-multi:symmetric-maps}
and~\ref{sec:higher-multi:tri-generalized} close the chapter with the
double-commutator theorem \cite[p.~187]{Vukman1990Two} and Wang
decomposition on triangular rings \cite[p.~1955]{Wang2016Bideriv}.
